% \iffalse meta-comment
%
% hit-report-heading.dtx 是开题报告的章节标题
%
% \fi
%
% \section{开题报告的章节标题}
%
% \changes{v3.2a}{2026/09/25}{深圳本科那两份报告的标题改由 \pkg{docxlike} 排，与终稿共用一层}
% \changes{v3.2a}{2026/09/13}{控制流改用 expl3（\cs{ctexset} 值里的条件由 ctex 延后求值，
%   保持原样）；威海本科中期跟着本部走；整块套 \option{stage} 闸；正文第一个编号标题
%   定顶级，从 \cs{chapter} 起的整体降一级，写乱或落到款以下报错}
%
% ^^A ======================================================================
% ^^A Module report-chapter: 章节标题格式（字号、缩进、间距、编号样式）（hit-report）
% ^^A ======================================================================
%    \begin{macrocode}
%<*report-chapter>
\bool_if:NF \g__hit_final_bool
  {
%    \end{macrocode}
%
%    \begin{macrocode}
\hit@if@shenzhen@doctor@interim@TF
  {
    \ctexset
      {
        section =
          {
            afterindent = true ,
            beforeskip = { 7.5bp } ,     % 上下空 0.5 行
            afterskip  = { 7.5bp } ,
            format = { \normalsize } ,
            aftername = \hit@after@number:nnn{section}{zh}{\enspace} ,
            fixskip = true ,
            break = { } ,
          } ,
        subsection =
          {
            afterindent = true ,
            beforeskip = { 7bp } , afterskip = { 7bp } ,
            format = { \normalsize } ,
            aftername = \hit@after@number:nnn{subsection}{zh}{\enspace} , fixskip = true , break = { } ,
          } ,
        subsubsection =
          {
            afterindent = true ,
            beforeskip = { 3bp } , afterskip = { 3bp } ,
            format = { \normalsize } ,
            aftername = \hit@after@number:nnn{subsubsection}{zh}{\enspace} , fixskip = true , break = { } ,
          } ,
        paragraph / afterindent = true ,
        subparagraph / afterindent = true
      }
  }
  {
    \dim_set:Nn \parindent { 2\ccwd }
    \ctexset
      {
        section =
          {
            afterindent = true ,
            beforeskip = { 7.5bp } ,     % 上下空 0.5 行
            afterskip  = { 7.5bp } ,
            format = { \heiti \hit@fontsize@x{xiaosan}{1.25} } ,
            aftername = \hit@after@number:nnn{section}{zh}{\enspace} ,
            fixskip = true ,
            break = { } ,
          } ,
        subsection =
          {
            afterindent = true ,
            beforeskip = { 7bp } , afterskip = { 7bp } ,
            format = { \heiti \hit@fontsize@x{sihao}{1.25} } ,
            aftername = \hit@after@number:nnn{subsection}{zh}{\enspace} , fixskip = true , break = { } ,
          } ,
        subsubsection =
          {
            afterindent = true ,
            beforeskip = { 3bp } , afterskip = { 3bp } ,
            format = { \heiti \normalsize } ,
            aftername = \hit@after@number:nnn{subsubsection}{zh}{\enspace} , fixskip = true , break = { } ,
          } ,
        paragraph / afterindent = true ,
        subparagraph / afterindent = true
      }
  }
%    \end{macrocode}
%
% \changes{v3.1d}{2025/03/03}{哈尔滨本科中期的一级节标题变成了中文……\issue[213]}
%    \begin{macrocode}
\bool_lazy_all:nT
  {
    { ! \bool_if_p:N \g__hit_shenzhen_bool }
    { \bool_if_p:N \g__hit_bachelor_bool }
    { \bool_if_p:N \g__hit_interim_bool }
  }
  {
    \ctexset
      { section = { number = \chinese { section } , name = { ,、} , aftername={\hit@after@number:nnn{section}{zh}{}} } }
  }
%    \end{macrocode}
%
% 深圳硕士的节标题格式。这一段的取值里有 \cs{ifhit@proposal} 与控制空格
% \cs{\ }，控制空格在 \cs{ExplSyntaxOn} 下的记号化与平时不同，实测会让
% 变体 32 的 PDF 产生差异。所以把整块取值放在 expl3 之外定义成一个宏，
% 记号与改动前逐个相同，再由 expl3 的条件来调用。
%    \begin{macrocode}
\cs_set:Npn \hit@report@shenzhen@master@section@settings
  {
%
~\ctexset{~
 section={~
 beforeskip=25bp,~
 afterskip=27pt,~
 format=\hit@if@option@TF{proposal}{\centering\xiaoer}{\xiaosan}\heiti,~
 aftername={\hit@after@number:nnn{section}{zh}{\hit@if@option@TF{proposal}{\quad}{\enspace}}},~
 break=\hit@if@option@TF{proposal}{\clearpage}{{}~}~
 },~
 subsection={~
 beforeskip=\hit@if@option@TF{proposal}{12pt~}{15pt~},~
 afterskip=\hit@if@option@TF{proposal}{12pt~}{15pt~},~
 format=\heiti\hit@if@option@TF{proposal}{\xiaosan}{\sihao},~
 fixskip=true,~
 aftername={\hit@after@number:nnn{subsection}{zh}{\hit@if@option@TF{proposal}{{\ }~}{\enspace}}}~
 },~
 subsubsection={~
 beforeskip=\hit@if@option@TF{proposal}{12pt~}{15pt~},~
 afterskip=\hit@if@option@TF{proposal}{12pt~}{15pt~},~
 format=\heiti\sihao,~
 fixskip=true,~
 aftername={\hit@after@number:nnn{subsubsection}{zh}{\hit@if@option@TF{proposal}{{\ }~}{\enspace}}}~
 }~
 }%
  }
\bool_lazy_and:nnT
  { \bool_if_p:N \g__hit_shenzhen_bool }
  { \bool_if_p:N \g__hit_master_bool }
  { \hit@report@shenzhen@master@section@settings }
%    \end{macrocode}
%
% 深圳博士中期在每个节标题上方画一条横线。第四个参数是星号，给了就不画，
% 原先在 \file{hithesisartplus.cls} 里。
%    \begin{macrocode}
\cs_new_protected:Npn \hit@report@shenzhen@doctor@section@settings
  {
    \cs_set_eq:NN \hit@section \section
    \RenewDocumentCommand \section { s o m s }
      {
        \IfBooleanF {##4}
          { \noindent \tikz [ overlay ] \draw ( -5mm , 0mm ) -- ( 155mm , 0mm ) ; }
        \IfBooleanTF {##1}
          {% \section*
            \IfNoValueTF {##2}
              { \hit@section * {##3} }
              { \hit@section * [##2] {##3} }
          }
          {% \section
            \IfNoValueTF {##2}
              { \hit@section {##3} }
              { \hit@section [##2] {##3} }
          }
      }
  }
\hit@if@shenzhen@doctor@interim { \hit@report@shenzhen@doctor@section@settings }
  }
%    \end{macrocode}
%
% \subsection{报告的标题层级}
%
% 报告没有章，顶级是节，往下是条、款。作者可以从 \cs{section} 写起，也可以从
% \cs{chapter} 写起：\emph{正文里第一个带编号的标题命令所在的那一级就是报告的
% 顶级}，其余顺延。从 \cs{chapter} 起的，\cs{chapter} 排成节、\cs{section} 排成条、
% \cs{subsection} 排成款；从 \cs{section} 起的照旧。这样一份从 \cs{chapter} 写起的
% 学位论文源码改个 \option{stage} 就能排成报告，老用户从 \cs{section} 写起的习惯也不变。
% 这条规则与 iota-hit 同一句话（那边按正文里编号标题的最小层级算偏移）。终稿不做这件事：
% 终稿必须有章，从 \cs{section} 写起的终稿不会被悄悄提成章。
%
% 偏移由第一个带编号的标题一次定死，不读 aux，编一遍就定。带星号的不编号标题不算，
% 它自己按“还没定”处理：\cs{chapter*} 排成节，其余照原级。
%
% 三种写法报错，规范没有规定的层级不给兜底：
% \begin{itemize}
% \item 从 \cs{section} 起头之后又写 \cs{chapter}：结构写乱了。出错后这个 \cs{chapter}
%   先当节排，后面的内容照常编下去，好看到别的错误。
% \item 第一个带编号的标题是 \cs{subsection} 或 \cs{subsubsection}：没有顶级。
% \item 降级之后落到款以下（从 \cs{chapter} 起头时的 \cs{subsubsection}）：规范只到款，
%   再往下写列表。\cs{paragraph} 与 \cs{subparagraph} 在两档里都一样，见
%   \file{src/engine/hit-loadclass.dtx}。
% \end{itemize}
%
% 类自己排的前置与后置（封面、目录、参考文献）不参与：\cs{hitfrontmatter} 与
% \cs{hitbackmatter} 排的时候按原级走，里面的 \cs{chapter*} 仍是 \pkg{ctex} 的章。
% 参考文献页的标题直接走顶级那一支，不经过偏移，正文里手写 \cs{printbibliography}
% 也不会被降成条。
%
% 正文照终稿排的那一档（深圳本科）顶级要另起一页，原先是 \cs{hitmainmatter} 往
% \cs{section} 前面补一次换页；现在 \cs{chapter} 也可能落到顶级，换页挪进顶级那一支，
% 由 \cs{g\_\_hit\_report\_top\_break\_bool} 开。
%    \begin{macrocode}
\int_new:N \g__hit_report_offset_int
\int_gset:Nn \g__hit_report_offset_int { -1 }
\bool_new:N \g__hit_report_heading_raw_bool
\bool_new:N \g__hit_report_top_break_bool
\cs_new_protected:Npn \__hit_report_heading_install:
  {
    \cs_gset_eq:NN \__hit_report_level_zero:w \chapter
    \cs_gset_eq:NN \__hit_report_level_one:w \section
    \cs_gset_eq:NN \__hit_report_level_two:w \subsection
    \cs_gset_eq:NN \__hit_report_level_three:w \subsubsection
    \cs_gset_protected:Npn \chapter
      { \__hit_report_heading:nn { 0 } { chapter } }
    \cs_gset_protected:Npn \section
      { \__hit_report_heading:nn { 1 } { section } }
    \cs_gset_protected:Npn \subsection
      { \__hit_report_heading:nn { 2 } { subsection } }
    \cs_gset_protected:Npn \subsubsection
      { \__hit_report_heading:nn { 3 } { subsubsection } }
  }
\cs_new_protected:Npn \__hit_report_heading:nn #1#2
  {
    \peek_charcode_remove:NTF *
      { \__hit_report_heading:nnn {#1} {#2} { * } }
      { \__hit_report_heading:nnn {#1} {#2} { } }
  }
\int_new:N \l__hit_report_level_int
\cs_new_protected:Npn \__hit_report_heading:nnn #1#2#3
  {
    \bool_if:NTF \g__hit_report_heading_raw_bool
      { \int_set:Nn \l__hit_report_level_int {#1} }
      {
        \tl_if_empty:nT {#3} { \__hit_report_offset_fix:nn {#1} {#2} }
        \int_compare:nNnTF { \g__hit_report_offset_int } < { 0 }
          {
            \int_set:Nn \l__hit_report_level_int
              { \int_max:nn {#1} { 1 } }
          }
          {
            \int_set:Nn \l__hit_report_level_int
              { \int_max:nn { #1 + \g__hit_report_offset_int } { 1 } }
          }
      }
    \__hit_report_heading_emit:nn {#2} {#3}
  }
\cs_new_protected:Npn \__hit_report_offset_fix:nn #1#2
  {
    \int_compare:nNnTF { \g__hit_report_offset_int } < { 0 }
      {
        \int_case:nnF {#1}
          {
            { 0 } { \int_gset:Nn \g__hit_report_offset_int { 1 } }
            { 1 } { \int_gset:Nn \g__hit_report_offset_int { 0 } }
          }
          {
            \msg_error:nnn { hithesis } { report-heading-no-top } {#2}
            \int_gset:Nn \g__hit_report_offset_int { 0 }
          }
      }
      {
        \bool_lazy_and:nnT
          { \int_compare_p:nNn {#1} = { 0 } }
          { \int_compare_p:nNn { \g__hit_report_offset_int } = { 0 } }
          { \msg_error:nn { hithesis } { report-chapter-after-section } }
      }
  }
\cs_new_protected:Npn \__hit_report_heading_emit:nn #1#2
  {
    \int_case:nnF { \l__hit_report_level_int }
      {
        { 0 } { \__hit_report_level_zero:w #2 }
        { 1 } { \hit@report@top@heading #2 }
        { 2 } { \__hit_report_level_two:w #2 }
        { 3 } { \__hit_report_level_three:w #2 }
      }
      {
        \msg_error:nnn { hithesis } { heading-too-deep } {#1}
        \hit@paragraph@original #2
      }
  }
\cs_new_protected:Npn \hit@report@top@heading
  {
    \bool_if:NT \g__hit_report_top_break_bool
      { \clearpage \bool_gset_true:N \g_docxlike_section_start_bool }
    \__hit_report_level_one:w
  }
%    \end{macrocode}
%
% \emph{正文照终稿排的那一档（深圳本科）的标题交给 \pkg{docxlike} 排}，共用的那一层在
% \file{hit-docxlike.dtx}，与终稿一样。节、条、款套“标题 1”到“标题 3”，编号来自一套三级的
% 列表“\%1”“\%1.\%2”“\%1.\%2.\%3”，列表第一级对 \texttt{section} 计数器。顶级另起一页之后立起
% “新一节开头”标记，标题在页顶留住段前（照 Word 减去前一段的段后），与终稿的章一样。其余几档的标题样式是空的，还是
% \pkg{ctex} 排，取值见上面那几段 \cs{ctexset}。
%    \begin{macrocode}
\cs_new_protected:Npn \__hit_report_heading_docxlike:
  {
    \clist_gset:Nn \g__hit_heading_levels_clist { section , subsection , subsubsection }
    \docxlikesetup
      {
        multilevel-lists = { hit-report = {
          level-1 = { enter-formatting-for-number = \% 1 , follow-number-with = nothing } ,
          level-2 = { enter-formatting-for-number = \% 1 . \% 2 , follow-number-with = nothing } ,
          level-3 = { enter-formatting-for-number = \% 1 . \% 2 . \% 3 , follow-number-with = nothing } ,
        } }
      }
    \__hit_heading_link:nn { hit-report } { 3 }
    \__hit_heading_styles:
    \__hit_heading_section_new:Nnnnn \section { section } { Heading~1 } { Heading1 } { 0 }
    \__hit_heading_section_new:Nnnnn \subsection { subsection } { Heading~2 } { Heading2 } { 1 }
    \__hit_heading_section_new:Nnnnn \subsubsection { subsubsection } { Heading~3 } { Heading3 } { 2 }
  }
\bool_if:NF \g__hit_final_bool
  {
    \bool_if:NT \g__hit_final_body_bool { \__hit_report_heading_docxlike: }
    \__hit_report_heading_install:
  }
%</report-chapter>
%    \end{macrocode}
%
% \Finale
